% ============================================================================
% DRAFT — Section "Data Collection and Methodology"  (Riccardo + Claude Code, 2026-09-29)
% Paste into Overleaf. Needs \usepackage{xcolor}. \TODO marks numbers/choices
% that are not final: fill them only once collection is stable.
% Everything described here matches what the scripts in the repo actually do
% (scripts/*.py, src/collect/*.py) as of 2026-09-29.
% ============================================================================
\providecommand{\TODO}[1]{\textcolor{red}{[TODO: #1]}}

\section{Data Collection and Methodology}
\label{sec:data}

Our goal is a corpus of \emph{paired observations}: a Telegram channel together
with the public group(s) it is associated with, both still active at collection
time. The pipeline combines automated steps, performed with a read-only
Telegram client, and manual steps, performed by the authors through the
official Telegram application. Table~\ref{tab:pipeline} \TODO{add table}
summarises which steps are automated and which are manual. Throughout the
collection the research account never joined a chat, sent a message, reacted,
clicked a link or requested access to a private chat.

\subsection{Seed Dataset}
\label{subsec:seed-dataset}

We use TGDataset~\cite{tgdataset} as a seed source. TGDataset contains
120{,}979 public Telegram channels and their messages, collected between
January 2021 and July 2022, and releases topic labels for the 19{,}768
English-language channels. It is the largest public collection of Telegram
channels with topic annotations, which lets us start from an externally defined
set of channels rather than from our own keyword choices.

TGDataset is a \emph{historical} seed: it does not describe Telegram in 2026,
and many of its channels have since been deleted, renamed or abandoned. We
therefore use it only to obtain candidate channels, and every candidate is
re-validated against the live platform (Section~\ref{subsec:channel-verification}).
In this pilot we consider the \emph{Crypto} topic (563 channels).
\TODO{other topics (e.g.\ Carding) if we extend beyond the pilot}

\subsection{Candidate Channel Selection}
\label{subsec:candidate-selection}

The topic labels identify channels only by numeric id, which is not sufficient
to open a channel through the Telegram API. We therefore stream the four
TGDataset archives from Zenodo (about 71\,GB compressed) and, for every channel
labelled \emph{Crypto}, extract only its metadata: username, title,
description, Telegram's \texttt{scam} and \texttt{verified} flags, number of
subscribers, date of its last message in the dataset, and the \texttt{t.me}
links that appear in its messages. No message content is retained at this stage.
This yields \TODO{N} candidates with a username \TODO{N without username, excluded}.

Because re-validation is rate-limited (Section~\ref{subsec:channel-verification}),
candidates are ranked offline before any request is made. The score favours
channels that were still posting in the last months covered by TGDataset, had
many subscribers, and linked to usernames that look like groups (e.g.\
containing \emph{chat}, \emph{group}, \emph{community}).
\TODO{decide: ranking only orders the queue, or also excludes low-score candidates}

\paragraph{Additional sources.}
TGDataset contains few channels that are both still active and associated with
a public group. We therefore add candidates from three live sources, recording
the source of every channel:
(i) Telegram's public search with \TODO{15} English crypto keywords
(e.g.\ \emph{crypto signals}, \emph{airdrop}, \emph{memecoin});
(ii) channels linked from the posts of the channels already verified, in
particular airdrop \emph{aggregators} that advertise many crypto projects, each
with its own channel and community group; and
(iii) channels that Telegram recommends as similar to a verified channel.
Duplicates are removed by username (case-insensitive).
\TODO{check: does the supervisor accept sources beyond TGDataset? (decisions.md, 2026-09-22)}

\subsection{Current Channel Verification}
\label{subsec:channel-verification}

Every candidate is verified automatically with a read-only client built on
Telethon~\cite{telethon}, using the MTProto API with a dedicated research
account. For each candidate the verifier:
\begin{enumerate}
  \item resolves the username; if it no longer exists, or now belongs to a user
        or bot, the channel is marked \emph{dead};
  \item checks that the username still points to the same channel id as in
        TGDataset (usernames can be reused after deletion);
  \item reads the public channel information (title, description,
        subscribers, flags, pinned message, linked discussion group);
  \item reads the last 200 posts and marks the channel \emph{active} if at
        least one was published in the 30 days before the verification date,
        \emph{inactive} otherwise.
\end{enumerate}
None of these operations requires joining the channel. The verification date
is stored with every record. \TODO{verification period: from ... to ...}

\paragraph{Rate limiting.}
All requests are serialised with a pause of 6\,s, and username resolutions
(the operation Telegram limits most strictly) are capped at 100 per day for
the whole account. During the pilot, about 250--300 resolutions in one day
triggered a 19-hour \texttt{FloodWait}; since then the collector stops at the
first wait longer than 60\,s and resumes the next day. Network errors are
recorded per channel; the run stops after three consecutive errors. All runs are
resumable and never re-verify a channel already verified.

\subsection{Group Discovery}
\label{subsec:group-discovery}

We call a group \emph{associated} with a channel when the channel explicitly
points to it and the group is public at collection time.
\TODO{confirm this definition with Marco/Simone; see decisions.md 2026-09-29}
For every active channel we search for such pointers in:
\begin{itemize}
  \item the channel description;
  \item the pinned message and the last 200 posts, including plain URLs, links
        hidden behind text, URL buttons and \texttt{@mentions};
  \item the channel's \emph{linked discussion group}, i.e.\ the group Telegram
        uses for comments.
\end{itemize}
Every username found is resolved and classified as group, channel, user, bot or
non-existent; private invite links (\texttt{t.me/+\ldots}) are recorded but never
opened. Discovery is performed automatically by the verifier and, in parallel,
manually by two authors who inspect each active channel in the Telegram app and
record the groups they find (Section~\ref{subsec:manual-validation}).

A linked discussion group is often only a comment section. We count it as a
group only if it has a public username and is used as a chat: at least
\TODO{10} messages that are not replies to channel posts, written by at least
\TODO{5} distinct users, in the last 7 days. A standalone group is considered
active with at least \TODO{20} messages from at least \TODO{5} distinct users in
the last 7 days. Activity is measured on the last 100 messages, which are
readable without joining.

\subsection{Manual Validation}
\label{subsec:manual-validation}

Two authors (Marco and Simone) inspect every active channel and, for each
candidate group, check that:
(i) the link resolves;
(ii) the destination is a group and not a channel or bot;
(iii) the group is public and readable;
(iv) the association with the source channel is explicit;
(v) the pair is not a duplicate.
They record the result, the date and their name in the shared master table.
The manual and automated results are compared pair by pair; disagreements are
resolved by discussion and reported as the error rate of the automated
discovery. \TODO{agreement measure (e.g.\ Cohen's $\kappa$ on a double-coded sample of N channels)}

\subsection{Dataset Construction}
\label{subsec:dataset-construction}

Figure~\ref{fig:funnel} \TODO{funnel figure} reports the number of channels
retained at each step:
candidate channels $\rightarrow$ existing channels $\rightarrow$ active
channels $\rightarrow$ channels pointing to a group $\rightarrow$ active public
groups $\rightarrow$ unique channel--group pairs.
\TODO{numbers once collection is stable}
% Preliminary, 2026-09-29 (NOT final, for us only):
%  - keyword search: 69 channels (42 active), 69 groups (8 standalone active);
%  - only 1 of ~43 active keyword channels linked a standalone public group;
%  - aggregators @AirdropDetective and @Airdrop: ~21 project channel + community pairs;
%  - first ~30 TGDataset crypto channels verified: mostly dead or inactive, 1 with an active group.
These counts describe how the corpus was built; they are not yet our research
results.

\subsection{Collected Information}
\label{subsec:collected-information}

For each chat we store only the following.
\emph{Channel and group metadata}: username, title, description, number of
subscribers/members, Telegram's \texttt{scam}/\texttt{fake}/\texttt{verified}
flags, linked chat, date of verification.
\emph{Messages} (last \TODO{90} days, one file per chat): message id, date,
sender type (user, the channel itself, anonymous admin), sender identifier,
admin signature when the channel shows it, text, reply target and thread,
forward source, views, forwards, number of replies, media \emph{type} (no media
file is downloaded), service-message type, and the URLs and \texttt{t.me} links
contained in the message.
\emph{Administrators}: the list of admins of each chat, where Telegram shows it
to non-members, with their admin title.
\emph{Direct messages}: unsolicited messages received by the research account,
with the same fields plus whether the sender is an admin or member of one of
the monitored chats; the account never replies.

We do not store user names, usernames, phone numbers, bios or profile
pictures. \TODO{pseudonymisation of user ids and redaction of phone numbers,
emails and user mentions in text: implemented, activation pending team decision}
